% !PRINTED 546
\weberpage{546}

\begin{center}
§. 165.

\gesperrt{Auflösung der cyklischen Gleichungen.}
\end{center}

Die \gesperrt{Lagrange}'schen Resolventen führen durch Anwendung der jetzt bewiesenen Sätze zu der Auflösung der cyklischen Gleichungen, genauer gesagt, zur Reduction auf \gesperrt{reine Gleichungen}.

Wir verstehen jetzt unter den $\alpha$ nicht mehr beliebige Variable, sondern die Wurzeln einer cyklischen Gleichung, so dass die cyklischen Functionen der $\alpha$ als bekannte Grössen zu betrachten sind.

Nach dem Theorem \S.\,164, 2. sind die Coëfficienten von $(\varepsilon,\alpha)^m$ cyklische Functionen der $\alpha$.

Verstehen wir unter $a_0, a_1\dots a_{m-1}$ Grössen in $\Omega$ und setzen
\weq{1}{\psi_\lambda = a_0 + a_1\varepsilon^\lambda + a_2\varepsilon^{2\lambda} + \cdots + a_{m-1}\varepsilon^{(m-1)\lambda},}
so folgt aus diesem Theorem
\weq{2}{(\varepsilon^\lambda,\alpha) = \sqrt[m]{\psi_\lambda}.}

Bezeichnet darin $\varepsilon$ eine primitive $m^{\mathrm{te}}$ Einheitswurzel, so sind in der Form $(\varepsilon^\lambda,\alpha)$ alle Resolventen enthalten.

Bemerken wir noch, dass $(1,\alpha) = a$ als die Summe der Wurzeln zu den bekannten Grössen gehört, so haben wir nach \S.\,164 (3)
\weq{3}{m\alpha = a + \sqrt[m]{\psi_1} + \sqrt[m]{\psi_2} + \cdots + \sqrt[m]{\psi_{m-1}},}
und damit also $\alpha$ durch Radicale $m^{\mathrm{ten}}$ Grades ausgedrückt, die unter den Wurzelzeichen ausser den Grössen, die von Hause aus in $\Omega$ vorkommen, noch $m^{\mathrm{te}}$ Einheitswurzeln enthalten.

Jedes dieser Radicale hat, für sich betrachtet, $m$ verschiedene Werthe, die sich um $m^{\mathrm{te}}$ Einheitswurzeln als Factoren von einander unterscheiden. Geben wir jeder $m^{\mathrm{ten}}$ Wurzel alle ihre Werthe, so erhalten wir aus (3) viele verschiedene Werthe von $\alpha$, unter denen nach \S.\,164, (4) die sämmtlichen Wurzeln $\alpha, \alpha_1, \alpha_2\dots\alpha_{m-1}$ vorkommen. Aber die Zahl der so aus (3) abgeleiteten Ausdrücke ist viel grösser, und es handelt sich noch darum, die beizubehaltenden von den auszusondernden zu unterscheiden. Am einfachsten führt dazu folgender Weg.
